\documentclass[10pt,a4paper]{amsart}
\usepackage[margin=19mm]{geometry}
\usepackage{amsmath,amssymb,mathtools}
\usepackage[scr=rsfs,cal=euler]{mathalfa}
\usepackage{xfrac}
\usepackage[unicode,hidelinks,bookmarks=false]{hyperref}
\newcommand{\ie}{i.e.\ }
\newcommand{\cf}{cf.\ }
\newcommand{\resp}{respectively\ }
\newcommand{\Subsection}{\subsection}
\providecommand{\nobreakdash}{\nobreak\hbox{-}}
\setlength{\emergencystretch}{2em}
\multlinegap0pt
\usepackage{bm}
\usepackage{bbm}
\usepackage{dsfont}
\usepackage{xspace}
\usepackage{cancel}
\usepackage{mathtools}
\usepackage{amsthm}
\makeatletter
\@ifundefined{theorem}{\newtheorem{theorem}{Theorem}}{}
\makeatother
\title[Some frustrating questions on dimensions of products of pose]{Some frustrating questions on dimensions of products of posets}
\author{George M. Bergman}
\date{Source: arXiv:2312.12615v4 (CC BY 4.0)}
\begin{document}
\maketitle

\noindent\emph{Source and licence.} Some frustrating questions on dimensions of products of posets. Version arXiv:2312.12615v4 (CC BY 4.0), distributed under the Creative Commons Attribution 4.0 International licence, as recorded in the arXiv licence metadata for this exact version and shown on the abstract page https://arxiv.org/abs/2312.12615v4. Full rights notice: licenses/arxiv-2312.12615-CC-BY-4.0.txt.\par\smallskip

\noindent\textit{Editorial scope.} Counted unit: the Theorem whose source label is \texttt{T.main} in the article (source0.tex lines 430--450, source label \texttt{T.main}), delivered together with the article's own complete original proof of it (source0.tex lines 452--564, one \texttt{proof} environment of 113 lines). No mathematics is written, repaired or re-derived: statement and proof are the article's own text line for line. Required context retained uncounted: none. Every definition, display and label of those lines is retained unchanged. Omitted from this package and not redistributed: 1--429, 451--451, 565--1804. Nothing inside a retained excerpt is omitted or altered: every retained body is delivered exactly as the article's lines are written, so no omission receipt of any kind accompanies this package. Citations: the retained excerpts carry no citation command at all, so no bibliography entry of the article is transcribed and no reference is redistributed. Reference adaptation: none was needed and none was made; every reference inside the delivered unit resolves inside it, so no unresolved reference or citation is produced. Printed slips kept literal: no printed word, symbol, hypothesis or number of the counted statement or of its proof was altered. Layout and class: the article's own class is \texttt{amsart} with packages including amsmath, amsfonts, amsthm, amssymb, indentfirst, epic, xurl, graphics, needspace, fontenc; the wrapper is the lane's pinned \texttt{amsart} editorial wrapper at 10pt on A4, which restates the theorem-like environments the retained text needs and adds the article's own macro definitions that the retained text needs (none), with extra wrapper packages (bm, bbm, dsfont, xspace, cancel, mathtools). The delivered numbering is the unit's own. Assets: no retained span contains a figure environment, a graphics-inclusion command, a PDF-page inclusion, a table or a TikZ picture; no image file is copied into this package, the package declares no asset dependency and no image file is redistributed with it. Licence: version arXiv:2312.12615v4 of this article is distributed under the Creative Commons Attribution 4.0 International licence, as recorded in the arXiv licence metadata for this exact version and shown on the abstract page https://arxiv.org/abs/2312.12615v4. A full rights notice is delivered at licenses/arxiv-2312.12615-CC-BY-4.0.txt. Characters: every retained excerpt is pure ASCII, so the delivered engine, XeLaTeX with its default Latin Modern text font, reports no missing character.
\begin{theorem}\label{T.main}
Let $d$ and $n$ be nonnegative integers, $(T_i)_{i\in d}$ a
$\!d\!$-tuple of chains, which for notational convenience we will
assume pairwise disjoint and $P$ a subposet of $\prod_{i\in d} T_i.$
Suppose there exists an $\!n\!$-tuple $(M_j)_{j\in n}$ of
pairwise disjoint subsets of $\bigcup_{i\in d} T_i$ such that
%
\begin{equation}\begin{minipage}[c]{35pc}\label{d.eq_coord}
For every comparable pair of elements $p\leq p'$
in $P,$ and every $j\in n,$ there exists $i\in d$ such
that the $\!i\!$-th coordinates of $p$ and $p'$ are equal,
and their common value is a member of $M_j.$
\end{minipage}\end{equation}

Then for any $\!n\!$-tuple of chains $(C_j)_{j\in n}$ we have
%
\begin{equation}\begin{minipage}[c]{35pc}\label{d.main}
$\dim(P\times\prod_{j\in n} C_j)\ \leq\ d.$
\end{minipage}\end{equation}
%
\end{theorem}

\begin{proof}
For notational simplicity, we may assume all the $C_j$ are equal
to a common chain $C,$ and that
$\bigcup_{j\in n} M_j = \bigcup_{i\in d} T_i.$
Indeed, given structures as in the statement of the theorem,
we may embed all the $C_j$ in a common chain $C,$ and if we
prove~\eqref{d.main} with the $C_j$ all replaced by $C,$
this will imply the same inequality for the original
choices of $C_j.$
Likewise, if we enlarge some of the $M_j$
(still keeping them disjoint) so that their
union becomes the whole set $\bigcup_{i\in d} T_i,$
then the hypothesis about pairs of elements $p<p'$ assumed
for the original choices of $M_j$ remains true.
So let us make those assumptions.

To obtain~\eqref{d.main}, we need an embedding of the poset
on the left-hand side, which we can now write
$P\times C^n,$ in a product of $d$ chains $T'_i$ $(i\in d).$
For this purpose we define
%
\begin{equation}\begin{minipage}[c]{35pc}\label{d.T'}
$T'_i\ =\ T_i\ltimes C,$ the set-theoretic product
of $T_i$ and $C,$ ordered lexicographically, i.e., so that
\end{minipage}\end{equation}
%
\begin{equation}\begin{minipage}[c]{35pc}\label{d.lex}
$(t,c)\ \leq\ (t',c')$ if and only if
$t<t',$ or $t=t'$ and $c\leq c'.$
\end{minipage}\end{equation}
%
Also, since the $M_j$ $(j\in n)$ partition
$\bigcup_{i\in d} T_i,$ we can set the notation
%
\begin{equation}\begin{minipage}[c]{35pc}\label{d.mu}
For each $t\in\bigcup_{i\in d} T_i,$ $m(t)\in n$ will denote the
unique value such that $t\in M_{m(t)}.$
\end{minipage}\end{equation}

We now define a map $f: P\times C^n \to \prod_{i\in d} T'_i$ as follows.
%
\begin{equation}\begin{minipage}[c]{35pc}\label{d.f}
For $p=(p_i)_{i\in d}\in P \subseteq\prod_{i\in d} T_i,$
and $c=(c_j)_{j\in n}\in C^n,$ let
$f(p,c)$ be the element of $\prod_{i\in d} T'_i$ whose
$\!i\!$-th coordinate is $(p_i,c_{m(p_i)})$ for each $i\in d.$
\end{minipage}\end{equation}
%
We wish to show that $f$ is an embedding of posets.

First, $f$ is isotone, i.e.,
%
\begin{equation}\begin{minipage}[c]{35pc}\label{d.f_isotone}
If $(p,c) \leq (p',c')$ in $P\times C^n,$ then
$f(p,c) \leq f(p',c')$ in $\prod_{i\in d} T'_i.$
\end{minipage}\end{equation}
%
This follows easily from~\eqref{d.lex} and~\eqref{d.f}.

To complete the proof that $f$ is an embedding
(and in particular, is one-to-one), we need to show that
%
\begin{equation}\begin{minipage}[c]{35pc}\label{d.notleq}
If $(p,c)\,\nleq\,(p',c'),$ \,then\, $f(p,c)\,\nleq\,f(p',c').$
\end{minipage}\end{equation}

This breaks down into two cases.
Suppose first that
%
\begin{equation}\begin{minipage}[c]{35pc}\label{d.p_notleq}
$p\ \nleq\ p'.$
\end{minipage}\end{equation}

In that case, for some $i$ we have $p_i\nleq p'_i,$
i.e., $p_i> p'_i,$ and
looking at the $\!i\!$-th coordinates of $f(p,c)$ and $f(p',c'),$
namely $(p_i,c_{m(p_i)})$ and $(p'_i,c'_{m(p'_i)}),$
we see from the lexicographic ordering of $T'_i$
that the former is $>$ the latter,
giving the conclusion of~\eqref{d.notleq}.

If, on the other hand,
%
\begin{equation}\begin{minipage}[c]{35pc}\label{d.p_leq}
$p\,\leq\,p',$ \,but\, $c\,\nleq\,c',$
\end{minipage}\end{equation}
%
then let us choose a $j\in n$ such that
%
\begin{equation}\begin{minipage}[c]{35pc}\label{d.c_j_notleq}
$c_j\,>\,c'_j.$
\end{minipage}\end{equation}
%
Now by the {\em first} condition of~\eqref{d.p_leq} and
the hypothesis~\eqref{d.eq_coord}, there is some $i\in d$ such that
%
\begin{equation}\begin{minipage}[c]{35pc}\label{d.use_eq_coord}
$p_i\,=\,p'_i\in M_j.$
\end{minipage}\end{equation}
%
Note that by~\eqref{d.f},
%
\begin{equation}\begin{minipage}[c]{35pc}\label{d.ith_coords}
$f(p,c)$ and $f(p',c')$ have $\!i\!$-th terms
$(p_i,c_j)$ and $(p'_i,c'_j)$ respectively.
\end{minipage}\end{equation}

Since the $\!T_i\!$-coordinates of the above two $\!i\!$-th terms
are the same by~\eqref{d.use_eq_coord}, the order-relation
between them is
that of the $\!C\!$-coordinates, which satisfy~\eqref{d.c_j_notleq}.
This completes the proof of~\eqref{d.notleq}, and of the Theorem.
\end{proof}

\end{document}
